% Sections 8.1 to 8.5, from lectures/L08/appendix/a_theory.md.

\section{Vad är en funktion?}
\label{c8:sec:funktion}
En \textbf{funktion} $f$ är en regel som för varje tillåtet värde på $x$ ger exakt ett värde
$f(x)$:
\[
  f : x \mapsto f(x)
\]
Funktioner kan representeras som algebraiska uttryck, tabeller eller grafer.

\begin{exempel}
$f(x) = 2x + 1$ ger $f(3) = 7$ och $f(-1) = -1$.
\end{exempel}

\section{Definitionsmängd och värdemängd}
\label{c8:sec:mangder}
\textbf{Definitionsmängden} $D_f$ är mängden av alla $x$-värden för vilka $f(x)$ är definierat.

\textbf{Värdemängden} $V_f$ är mängden av alla möjliga funktionsvärden $f(x)$.

\begin{narrowtable}{@{}p{7em}p{8.5em}p{8.5em}@{}}
  \thead{Funktion} & \thead{Definitionsmängd} & \thead{Värdemängd} \\
  \midrule
  $f(x) = 2x + 3$ & $\mathbb{R}$ & $\mathbb{R}$ \\
  $f(x) = x^2$ & $\mathbb{R}$ & $[0, \infty)$ \\
  $f(x) = \sqrt{x}$ & $[0, \infty)$ & $[0, \infty)$ \\
  $f(x) = \dfrac{1}{x}$ & $\mathbb{R} \setminus \{0\}$ & $\mathbb{R} \setminus \{0\}$ \\
\end{narrowtable}

\section{Vanliga funktionstyper}
\label{c8:sec:funktionstyper}

\subsection{Linjär funktion}
\label{c8:sec:linjar}
\[
  f(x) = kx + m
\]
där $k$ är lutningen och $m$ är skärningspunkten med $y$-axeln.

\subsection{Potensfunktion}
\label{c8:sec:potens}
\[
  f(x) = x^n
\]

\subsection{Exponentialfunktion}
\label{c8:sec:exponential}
\[
  f(x) = C \cdot a^x, \qquad a > 0,\ a \neq 1
\]
\begin{itemize}
  \item $a > 1$: växande funktion.
  \item $0 < a < 1$: avtagande funktion.
  \item $f(0) = C$, startvärdet.
\end{itemize}

\subsection{Polynomfunktion}
\label{c8:sec:polynom}
\[
  f(x) = a_n x^n + a_{n-1} x^{n-1} + \cdots + a_0
\]

\section{Att läsa en graf}
\label{c8:sec:graf}
Från en graf avläser man:
\begin{itemize}
  \item definitionsmängden, alltså vilka $x$-värden som finns,
  \item värdemängden, alltså vilka $y$-värden som uppnås,
  \item nollställena, de $x$ där $f(x) = 0$,
  \item maximum- och minimumpunkterna.
\end{itemize}

\section{Typexempel}
\label{c8:sec:typexempel}

\begin{exempel}[Temperatursensor]
En sensor ger temperaturen $f(x)$ (i $^{\circ}$C) som funktion av inspänningen $x$ (i V):
\[
  f(x) = 100x - 50, \qquad 0 \leq x \leq 5\,\text{V}
\]
\textbf{a)} Bestäm definitionsmängd och värdemängd. Definitionsmängden är
$D_f = [0, 5]\,\text{V}$, och
\[
  f(0) = -50\,^{\circ}\text{C}, \quad f(5) = 450\,^{\circ}\text{C}
  \quad \Rightarrow \quad V_f = [-50, 450]\,^{\circ}\text{C}.
\]
\textbf{b)} Beräkna inspänningen när temperaturen är $30\,^{\circ}\text{C}$.
\[
  100x - 50 = 30 \quad \Rightarrow \quad x = 0{,}8\,\text{V}
\]
\end{exempel}

\begin{exempel}[Linjär funktion från två punkter]
En linjär funktion passerar genom $(1, 3)$ och $(3, 7)$. Lutningen är
\[
  k = \frac{7 - 3}{3 - 1} = 2.
\]
Insättning av $(1, 3)$ ger $m = 3 - 2 \cdot 1 = 1$, alltså
\[
  f(x) = 2x + 1.
\]
\end{exempel}

\begin{exempel}[Exponentialfunktion och tillväxt]
En datamängd startar med $6 \times 10^6$ datapunkter och växer med $3\,\%$ per år:
\[
  f(x) = 6 \times 10^6 \cdot 1{,}03^x
\]
\textbf{a)} Beräkna $f(5)$.
\[
  f(5) = 6 \times 10^6 \cdot 1{,}03^5 \approx 6{,}96 \times 10^6
\]
\textbf{b)} När har datamängden fördubblats?
\[
  1{,}03^x = 2 \quad \Rightarrow \quad x = \frac{\log 2}{\log 1{,}03} \approx 23{,}5\,\text{år}
\]
\end{exempel}
